\SourceLeaf{471}
\par\Needspace{12\baselineskip}\bigskip
\pdfbookmark[2]{Darstellung, Vierter Abschnitt, Erster Entwurf}{pallas-II-section4-first}
\begin{center}
{\large\bfseries Vierter Abschnitt.\quad [Erster Entwurf.]\par}
\medskip
\textit{Grundsätze zur Bestimmung der Bewegung für unbegrenzte Zeit.}
\end{center}
\pdfbookmark[3]{Darstellung, [21]}{pallas-II-article21}
\begin{center}[21.]\end{center}
\noindent\hspace*{\parindent}Die Integration der momentanen Änderungen der Elemente anhand ihrer Zahlenwerte liefert die Elemente nur für eine begrenzte Zeit. Um die unbestimmte Integration ausführen zu können, muss man diese Änderungen zunächst als integrierbare Funktionen der Zeit ausdrücken. Darin liegt die eigentliche Schwierigkeit der Aufgabe. Angesichts der Kleinheit der störenden Massen im Vergleich zur Masse der Sonne betrachten wir die Änderungen der Elemente oder allgemein die Größen, die den Faktor \(\mu\) enthalten, als sehr kleine Größen \textit{erster Ordnung}. Größen, die sein Quadrat oder das Produkt von \(\mu\) mit einer anderen störenden Masse enthalten, sind von zweiter Ordnung und so fort. Von nun an bezeichnen wir den wahren Wert der Länge des Knotens (auf der festen Ekliptik gezählt) für einen unbestimmten Zeitpunkt mit \(\AscendingNode+\delta\AscendingNode\). Dabei ist \(\AscendingNode\) eine konstante Größe und \(\delta\AscendingNode\) der veränderliche Teil erster Ordnung. Bei den anderen Elementen verfahren wir ebenso. Daraus ist ersichtlich, dass, wenn man die momentane Änderung eines Elements bestimmt, indem man alle in ihren Wert eingehenden Größen aus den konstanten Teilen der Elemente berechnet, diese Änderung bis zur zweiten Ordnung (ausschließlich) richtig ist. Allgemein liegt der Fehler der Änderungen der Elemente um eine Ordnung höher als die Ordnung der bei ihrer Berechnung verwendeten Werte der Elemente. Daraus ergibt sich die Methode der sukzessiven Näherungen. Man muss jedoch einräumen, dass die Ausführung dieser im Ansatz so einfachen Methode fast undurchführbar werden könnte, wenn man auch nur die zweite Näherung \textit{vollständig} berechnen müsste. Für die Bedürfnisse der Astronomie ist dies glücklicherweise keineswegs notwendig. Man beschränkt sich auf die erste Näherung und fügt einige Abänderungen hinzu, die sich aus einer begründeten Betrachtung der Glieder zweiter Ordnung ergeben.
